% 2021 全国甲卷理科 2：农村家庭年收入频率分布直方图
% 根据固定上游 PDF 第 1 页：组距均为 1，组界 2.5 至 14.5。
% 各柱高依次为 0.02, 0.04, 0.10, 0.14, 0.20, 0.20, 0.10, 0.10, 0.04, 0.02, 0.02, 0.02。
\begin{tikzpicture}[font=\small]
  \def\xs{0.48}\def\ys{15.0}\def\xz{0.52}
  % 横轴断轴与纵轴
  \draw (0,0) -- (0.10,0);
  \draw (0.10,0) -- (0.16,0.09) -- (0.25,-0.09) -- (0.31,0);
  \draw[->] (0.31,0) -- ({\xz+12*\xs+0.55},0);
  \draw[->] (0,0) -- (0,{0.20*\ys+0.45});
  \node[anchor=north] at (0,-0.04) {$0$};
  \node[anchor=south east,align=center] at (-0.04,{0.20*\ys+0.1}) {频率/组距};
  \node[anchor=west] at ({\xz+12*\xs+0.2},0.15) {收入/万元};
  % 刻度参考线及标注
  \foreach \v in {0.02,0.04,0.10,0.14,0.20} {
    \draw[densely dotted,gray!70] (0,{\v*\ys}) -- ({\xz+12*\xs},{\v*\ys});
    \node[anchor=east,inner sep=2pt,font=\scriptsize] at (-0.04,{\v*\ys}) {$\v$};
  }
  % 十二组柱形
  \foreach \i/\h in {0/0.02,1/0.04,2/0.10,3/0.14,4/0.20,5/0.20,
                       6/0.10,7/0.10,8/0.04,9/0.02,10/0.02,11/0.02} {
    \draw ({\xz+\i*\xs},0) rectangle ({\xz+(\i+1)*\xs},{\h*\ys});
  }
  \foreach \i/\xlab in {0/2.5,1/3.5,2/4.5,3/5.5,4/6.5,5/7.5,6/8.5,
                         7/9.5,8/10.5,9/11.5,10/12.5,11/13.5,12/14.5} {
    \node[anchor=north,inner sep=1pt,font=\tiny] at ({\xz+\i*\xs},-0.04) {$\xlab$};
  }
\end{tikzpicture}
